\section{Задание 4. Ядро и образ линейного оператора}

\textbf{Задача.}

Линейный оператор $\varphi$ задан в стандартном базисе пространства $\mathbb{R}^5$ матрицей $A$.

1. Найти ядро линейного оператора, указать его размерность.

2. Найти образ линейного оператора, указать его размерность.

3. Убедиться, что выполняется теорема о размерностях ядра и образа линейного оператора.

\[
17.\quad
A=
\begin{pmatrix}
0 & 1 & -1 & 1 & 1\\
1 & -1 & 0 & 2 & -1\\
-2 & 2 & 1 & 5 & 0\\
-4 & 4 & 1 & 1 & 2\\
-2 & 0 & 3 & 3 & -2
\end{pmatrix}.
\]

\textbf{Решение.}

\textbf{1. Найдём ядро линейного оператора.}

Найдем сначала ядро оператора $A$. Согласно определению
\[
\ker\varphi=\{x\in\mathbb{R}^5:\varphi(x)=0\},
\]
то есть для нахождения ядра оператора $A$ необходимо найти множество решений системы
\[
Ax=0.
\]

Для этого с помощью элементарных преобразований строк приведем матрицу $A$ к ступенчатому виду и найдем ее ранг:
\[
A=
\begin{pmatrix}
0 & 1 & -1 & 1 & 1\\
1 & -1 & 0 & 2 & -1\\
-2 & 2 & 1 & 5 & 0\\
-4 & 4 & 1 & 1 & 2\\
-2 & 0 & 3 & 3 & -2
\end{pmatrix}
\sim
\begin{pmatrix}
1 & -1 & 0 & 2 & -1\\
0 & 1 & -1 & 1 & 1\\
-2 & 2 & 1 & 5 & 0\\
-4 & 4 & 1 & 1 & 2\\
-2 & 0 & 3 & 3 & -2
\end{pmatrix}
\]
\[
\sim
\begin{pmatrix}
1 & -1 & 0 & 2 & -1\\
0 & 1 & -1 & 1 & 1\\
0 & 0 & 1 & 9 & -2\\
0 & 0 & 1 & 9 & -2\\
0 & -2 & 3 & 7 & -4
\end{pmatrix}
\sim
\begin{pmatrix}
1 & -1 & 0 & 2 & -1\\
0 & 1 & -1 & 1 & 1\\
0 & 0 & 1 & 9 & -2\\
0 & 0 & 1 & 9 & -2\\
0 & 0 & 1 & 9 & -2
\end{pmatrix}
\]
\[
\sim
\begin{pmatrix}
1 & -1 & 0 & 2 & -1\\
0 & 1 & -1 & 1 & 1\\
0 & 0 & 1 & 9 & -2\\
0 & 0 & 0 & 0 & 0\\
0 & 0 & 0 & 0 & 0
\end{pmatrix}.
\]

Ранг матрицы $A$ равен $3$, следовательно,
\[
\dim\ker\varphi=5-3=2.
\]

Пространство решений системы $Ax=0$ задается системой
\[
\left\{
\begin{aligned}
x_1-x_2+2x_4-x_5&=0,\\
x_2-x_3+x_4+x_5&=0,\\
x_3+9x_4-2x_5&=0.
\end{aligned}
\right.
\]

Выразим главные переменные через свободные $x_4$ и $x_5$:
\[
\left\{
\begin{aligned}
x_3&=-9x_4+2x_5,\\
x_2&=x_3-x_4-x_5=-10x_4+x_5,\\
x_1&=x_2-2x_4+x_5=-12x_4+2x_5.
\end{aligned}
\right.
\]

Тогда
\[
x=
x_4
\begin{pmatrix}
-12\\
-10\\
-9\\
1\\
0
\end{pmatrix}
+
x_5
\begin{pmatrix}
2\\
1\\
2\\
0\\
1
\end{pmatrix}.
\]

Отсюда получаем базис ядра:
\[
a_1=
\begin{pmatrix}
-12\\
-10\\
-9\\
1\\
0
\end{pmatrix},
\qquad
a_2=
\begin{pmatrix}
2\\
1\\
2\\
0\\
1
\end{pmatrix}.
\]

Следовательно,
\[
\ker\varphi=
\left\langle
\begin{pmatrix}
-12\\
-10\\
-9\\
1\\
0
\end{pmatrix},
\begin{pmatrix}
2\\
1\\
2\\
0\\
1
\end{pmatrix}
\right\rangle.
\]

\bigskip

\textbf{2. Найдём образ линейного оператора.}

Теперь найдем образ оператора $A$. Образ оператора порожден столбцами матрицы $A$.
Так как в ступенчатой матрице главные элементы стоят в первом, втором и третьем столбцах, то соответствующие столбцы исходной матрицы $A$ образуют базис образа.

Следовательно,
\[
b_1=
\begin{pmatrix}
0\\
1\\
-2\\
-4\\
-2
\end{pmatrix},
\qquad
b_2=
\begin{pmatrix}
1\\
-1\\
2\\
4\\
0
\end{pmatrix},
\qquad
b_3=
\begin{pmatrix}
-1\\
0\\
1\\
1\\
3
\end{pmatrix}.
\]

Таким образом,
\[
\Im\varphi=
\left\langle
\begin{pmatrix}
0\\
1\\
-2\\
-4\\
-2
\end{pmatrix},
\begin{pmatrix}
1\\
-1\\
2\\
4\\
0
\end{pmatrix},
\begin{pmatrix}
-1\\
0\\
1\\
1\\
3
\end{pmatrix}
\right\rangle,
\qquad
\dim\Im\varphi=3.
\]

\bigskip

\textbf{3. Убедимся, что выполняется теорема о размерностях ядра и образа линейного оператора.}

Проверим теорему о размерностях ядра и образа линейного оператора:
\[
\dim\ker\varphi+\dim\Im\varphi=2+3=5.
\]

Так как
\[
2+3=5=\dim\mathbb{R}^5,
\]
теорема о размерностях выполняется.
